% =====================================================================
\chapter{Embeddings, analogías y atención}
\label{tema:5}
% =====================================================================

Este tema aplica las herramientas anteriores a la entrada y la salida de un modelo de lenguaje. La capa de embeddings es una matriz; leer el vector de un token es multiplicar por un vector \emph{one-hot}; las relaciones semánticas aparecen, de forma aproximada, como direcciones del espacio; y la atención compara tokens mediante productos escalares. La presentación sigue el capítulo~5 de Jurafsky y Martin~\cite{jm3} y la clase~2 de CS224N~\cite{cs224n}, junto con los artículos de Mikolov y coautores~\cite{mikolov-a,mikolov-b} y de Vaswani y coautores~\cite{vaswani}.

\section{La matriz de embeddings}

En este tema se usa la notación del módulo~1: $\V$ es el vocabulario, $K=\abs{\V}$ su tamaño y los tokens se identifican con sus índices $0,1,\dots,K-1$. Las filas de las matrices indexadas por tokens se numeran también desde~0.

\begin{definicion}[Matriz de embeddings]\label{def:embedding}
Sea $d\ge 1$. Una \emph{matriz de embeddings} es una matriz $E\in\R^{K\times d}$. El \emph{embedding} del token $i$ es la fila $i$ de $E$, que se denota $E_{i,:}\in\R^{d}$ (escrita como columna cuando se opera con ella).
\end{definicion}

\begin{definicion}[Vector one-hot]
El vector \emph{one-hot} del token $i$ es el vector $e_i\in\R^{K}$ que tiene un~1 en la posición $i$ y ceros en el resto.
\end{definicion}

\begin{proposicion}\label{prop:onehot}
Para todo token $i$, $e_i^{\top}E=E_{i,:}^{\top}$; es decir, multiplicar el vector one-hot de un token por la matriz de embeddings selecciona su fila. Si $x_{1:n}$ es una secuencia de tokens y $O\in\R^{n\times K}$ es la matriz cuyas filas son $e_{x_1}^{\top},\dots,e_{x_n}^{\top}$, entonces $X=OE\in\R^{n\times d}$ es la matriz cuya fila $t$ es el embedding de $x_t$.
\end{proposicion}

\begin{proof}
La componente $j$ de $e_i^{\top}E$ es $\sum_k(e_i)_kE_{kj}$. Todos los sumandos con $k\neq i$ se anulan y queda $E_{ij}$. La segunda afirmación se sigue de la \cref{prop:cuatro}(ii): la fila $t$ de $OE$ es $e_{x_t}^{\top}E$.
\end{proof}

\begin{ejemplo}\label{ej:embedding}
Sean $K=4$, $d=3$ y
\[
E=\begin{pmatrix}0{,}1&0{,}2&0\\1&0&-1\\0{,}5&0{,}5&0{,}5\\-0{,}2&0{,}3&0{,}1\end{pmatrix}.
\]
El token~2 tiene vector one-hot $e_2=(0,0,1,0)$ y $e_2^{\top}E=(0{,}5;\,0{,}5;\,0{,}5)$. La secuencia $(2,0,2)$ produce una matriz $X\in\R^{3\times 3}$ cuyas filas primera y tercera coinciden: sin información de posición, el modelo no distingue las dos apariciones del token~2 (módulo~14).
\end{ejemplo}

\begin{observacion}[Indexar en lugar de multiplicar]
El producto $e_i^{\top}E$ requiere $Kd$ multiplicaciones, casi todas por cero; leer directamente la fila $i$ solo copia $d$ números. Las bibliotecas de redes neuronales implementan la capa de embeddings como una indexación, aunque matemáticamente sea el producto de la \cref{prop:onehot}. En GPT-2, con $K=50\,257$~\cite[§2.3]{radford} y $d=768$, la matriz $E$ tiene $50\,257\cdot 768=38\,597\,376$ parámetros, unos $154$~MB en precisión de 32 bits.
\end{observacion}

\section{De dónde salen los embeddings}

Los embeddings no se diseñan a mano como en el \cref{ej:juguete}: se inicializan al azar y se ajustan por descenso de gradiente (módulos~4 y~7) para que el modelo prediga bien el token siguiente. La justificación de que este procedimiento produzca vectores con significado es la \emph{hipótesis distribucional}: las palabras que aparecen en contextos parecidos tienden a tener significados parecidos~\cite{harris}. Firth la resumió en la frase «you shall know a word by the company it keeps»~\cite{firth}. Dos tokens que aparecen en contextos similares reciben gradientes similares y acaban con embeddings próximos en el sentido del coseno~\cite[cap.~5]{jm3}.

Antes de los Transformers, los embeddings se entrenaban por separado con modelos como word2vec~\cite{mikolov-a}, que aprende un vector por palabra prediciendo las palabras vecinas en una ventana. En un LLM, en cambio, la matriz $E$ se aprende junto con el resto de los parámetros.

\section{La capa de salida y el atado de pesos}

Al final del modelo, el vector $h\in\R^{d}$ de la última posición se compara con un vector por cada token del vocabulario. Si $U\in\R^{K\times d}$ es la matriz de salida, los logits son
\[
z=Uh\in\R^{K},\qquad z_i=U_{i,:}\cdot h,
\]
y la softmax del módulo~1 (tema~4) los convierte en una distribución sobre el token siguiente. Cada logit es, por tanto, un producto escalar entre $h$ y un vector asociado al token.

\begin{definicion}[Atado de pesos]
Se dice que un modelo \emph{ata} los pesos de entrada y salida (\emph{weight tying}) si usa $U=E$: la misma matriz sirve para leer los embeddings de entrada y para calcular los logits.
\end{definicion}

Press y Wolf~\cite{press} mostraron que esta elección mejora los modelos de lenguaje recurrentes, y el Transformer original la adopta~\cite[§3.4]{vaswani}. Ahorra $Kd$ parámetros, que en GPT-2 son $38\,597\,376$, y da a ambos usos una interpretación común: el logit del token $i$ es alto cuando $h$ apunta en la dirección del embedding de $i$.

\section{Analogías y aritmética de vectores}

Sea $v_w$ el embedding de la palabra $w$. Si un rasgo como el género corresponde aproximadamente a una dirección del espacio, la diferencia $v_{\text{mujer}}-v_{\text{hombre}}$ debería parecerse a $v_{\text{reina}}-v_{\text{rey}}$. Mikolov, Yih y Zweig~\cite{mikolov-b} observaron este fenómeno en embeddings aprendidos y propusieron resolver analogías con el método siguiente.

\begin{definicion}[Analogía por desplazamiento de vectores]\label{def:analogia}
Dada la pregunta «$a$ es a $b$ como $?$ es a $c$», se calcula $q=v_a-v_b+v_c$ y se responde con
\[
\hat w=\operatorname*{arg\,max}_{w\in\V\setminus\{a,b,c\}}\cos(v_w,q).
\]
\end{definicion}

\begin{proposicion}\label{prop:paralelogramo}
Si $v_a-v_b=v_{w^{*}}-v_c$, entonces $q=v_{w^{*}}$ y $\cos(v_{w^{*}},q)=1$, de modo que $w^{*}$ es una respuesta de máxima similitud.
\end{proposicion}

\begin{proof}
$q=v_a-v_b+v_c=(v_{w^{*}}-v_c)+v_c=v_{w^{*}}$, y $\cos(u,u)=1$ para $u\neq 0$; por el \cref{teo:cs}, ningún candidato supera este valor.
\end{proof}

\begin{ejemplo}
Con los vectores del \cref{ej:juguete}, $v_{\text{rey}}-v_{\text{hombre}}+v_{\text{mujer}}=(2,1,2,0)-(0,1,2,0)+(0,-1,2,0)=(2,-1,2,0)=v_{\text{reina}}$. La resta elimina «masculino» y «persona»; la suma añade «femenino» y «persona»; la realeza se conserva.
\end{ejemplo}

\begin{observacion}[Limitaciones]
Con embeddings reales, la igualdad de la \cref{prop:paralelogramo} solo se cumple de forma aproximada, y $q$ suele quedar más cerca de $v_a$ que de la respuesta esperada; la exclusión de $a$, $b$ y $c$ en la \cref{def:analogia} es esencial para que el método «funcione». Nissim, van Noord y van der Goot~\cite{nissim} analizan cómo esta exclusión distorsiona las conclusiones que se extraen de las analogías. En el cuaderno de prácticas del módulo, con 40 palabras en 16 dimensiones y ruido, la consulta rey $-$ hombre $+$ mujer da «reina» con coseno $0{,}924$, pero sin la exclusión «rey» queda segunda con $0{,}835$. Además, los embeddings heredan los sesgos de los textos de entrenamiento: Bolukbasi y coautores~\cite{bolukbasi} encontraron analogías que reproducen estereotipos de género. La idea de que los conceptos son direcciones reaparece en la interpretabilidad de los LLM (módulo~23).
\end{observacion}

\section{Adelanto: el producto escalar en la atención}

La atención (módulo~13) es un mecanismo de búsqueda por similitud. Cada token emite una consulta $q\in\R^{d_k}$, la compara por producto escalar con las claves $k_1,\dots,k_n$ de los tokens anteriores y reparte su atención con una softmax:
\[
s_j=\frac{q\cdot k_j}{\sqrt{d_k}},\qquad \alpha=\softmax(s_1,\dots,s_n).
\]
Para todos los tokens a la vez, las puntuaciones forman la matriz $QK^{\top}/\sqrt{d_k}$ (\cref{def:gram}).

\begin{ejemplo}\label{ej:atencion}
Sean $d_k=2$, $q=(1,0)$ y las claves $k_1=(1,1)$, $k_2=(0,2)$, $k_3=(-1,0)$. Los productos escalares son $1$, $0$ y $-1$; divididos por $\sqrt2$ dan $(0{,}7071;\,0;\,-0{,}7071)$, y la softmax da $\alpha\approx(0{,}5760;\,0{,}2840;\,0{,}1400)$. El token cuya clave está más alineada con la consulta recibe más atención.
\end{ejemplo}

El factor $1/\sqrt{d_k}$ tiene una justificación estadística~\cite[§3.2.1]{vaswani}: si las componentes de $q$ y $k$ son variables independientes de media~0 y varianza~1, entonces $\Var(q\cdot k)=d_k$ (ejercicio~\ref{eje:5-atencion}). Sin el factor, para $d_k$ grande las puntuaciones tendrían valores muy dispares, la softmax quedaría casi concentrada en un solo token y los gradientes serían casi nulos. Dividir por $\sqrt{d_k}$ devuelve la varianza a~1.

\begin{lecturas}
\item Jurafsky y Martin~\cite{jm3}, cap.~5 «Embeddings».
\item CS224N~\cite{cs224n}, clase~2 «Word Vectors» y tarea~1.
\item Zhang \emph{et al.}~\cite{d2l}, §15.7 «Word Similarity and Analogy».
\item Mikolov, Yih y Zweig~\cite{mikolov-b}; Vaswani \emph{et al.}~\cite{vaswani}, §3.2.1 y §3.4.
\end{lecturas}

\section*{Ejercicios del tema 5}
\addcontentsline{toc}{section}{Ejercicios del tema 5}

\begin{ejercicio}{1}{eje:5-tamano}
GPT-2 tiene $K=50\,257$ y, en su versión más pequeña, $d=768$. ¿Qué tamaño tiene $E$? ¿Cuántos parámetros tiene y cuánto ocupa con 4 bytes por número? Repita el cálculo con $d=1600$, la dimensión de la versión más grande.
\end{ejercicio}

\begin{ejercicio}{1}{eje:5-onehot}
Explique, a partir de la definición de producto de matrices, por qué $e_i^{\top}E$ es la fila $i$ de $E$, y estime cuántas veces más operaciones requiere el producto que la lectura directa de la fila en GPT-2.
\end{ejercicio}

\begin{ejercicio}{1}{eje:5-busqueda}
Se desea encontrar las palabras más similares a una consulta $q$ en un vocabulario muy grande. (a) Explique por qué, si se normalizan una vez todas las filas de $E$ para formar $\hat E$, los $K$ cosenos se obtienen con un único producto matriz-vector. (b) Cuente las multiplicaciones necesarias con $K=50\,257$ y $d=768$. (c) ¿Cómo se procesa a la vez un lote de $B$ consultas?
\fuente{adaptado de Zhang \emph{et al.}~\cite{d2l}, §15.7, ejercicio~2.}
\end{ejercicio}

\begin{ejercicio}{2}{eje:5-analogias}
Con los vectores del \cref{ej:juguete}: (a) calcule $\cos(\text{hombre},\text{mujer})$; (b) compruebe que $v_{\text{rey}}-v_{\text{hombre}}+v_{\text{mujer}}=v_{\text{reina}}$. (c) Con $v_{\text{España}}=(1,0,1,0)$, $v_{\text{Madrid}}=(0,1,1,0)$, $v_{\text{Francia}}=(1,0,0,1)$, $v_{\text{París}}=(0,1,0,1)$ y $v_{\text{Lyon}}=(0;\,0{,}5;\,0;\,1)$, calcule $q=v_{\text{Madrid}}-v_{\text{España}}+v_{\text{Francia}}$ y su coseno con París y con Lyon.
\end{ejercicio}

\begin{ejercicio}{2}{eje:5-secuencia}
Con la matriz $E$ del \cref{ej:embedding}: (a) calcule $e_2^{\top}E$ columna a columna; (b) escriba la matriz $O$ de la \cref{prop:onehot} para la secuencia $(2,0,2)$ y calcule $X=OE$ indicando su tamaño.
\end{ejercicio}

\begin{ejercicio}{2}{eje:5-logits}
Un modelo con pesos atados usa la matriz $E$ del \cref{ej:embedding} como matriz de salida. El vector final es $h=(1,0,-1)$. Calcule los logits $z=Eh$ y la distribución $\softmax(z)$. ¿Qué token es el más probable y por qué, en términos de ángulos?
\end{ejercicio}

\begin{ejercicio}{3}{eje:5-atencion}
(a) Reproduzca los cálculos del \cref{ej:atencion}. (b) Sean $q,k\in\R^{d}$ vectores aleatorios cuyas $2d$ componentes son independientes, de media~0 y varianza~1. Demuestre que $\E[q\cdot k]=0$ y $\Var(q\cdot k)=d$. (c) Con $d=64$, compare la softmax de las puntuaciones típicas $(8,0,-8)$, del orden de una desviación típica, con la de $(1,0,-1)$, que se obtiene al dividir por $\sqrt{64}$.
\end{ejercicio}
